% Generated by `taskgen course examples` from lattices.babai_easy (seed 1).
% Do not edit: change examples.yaml and regenerate.
\begin{example}[алгоритм Бабаи, хороший базис]\label{ex:lattice-babai-good}
В решётке $\mathcal{L}$, заданной базисными векторами \[v_1 = \begin{pmatrix}0 \\ 61\end{pmatrix}, v_2 = \begin{pmatrix}-97 \\ -5\end{pmatrix},\] найти вектор, ближайший к вектору $w = (-21, 99)^T$.

\par\noindent\textit{Решение.} $\det \mathcal{L} = 5917$, $HR(\mathcal{L}) = \left(\frac{|\det\mathcal{L}|}{\|v_1\|\|v_2\|}\right)^{1/2} = 0.999$ --- базис хороший (почти ортогональный). Применяем алгоритм Бабаи к исходному базису $u_i = v_i$. Решаем $w = t_1 u_1 + t_2 u_2$ (правило Крамера): $t_1 = \frac{9708}{5917} \approx 1.641$, $t_2 = \frac{21}{97} \approx 0.216$; округляем: $2 u_1 + 0 u_2 = (0, 122)^T$, $\|w - r\| \approx 31.145$.\[r = (0, 122)^T,\quad HR(\mathcal{L}) = 0.999.\]
\end{example}
